
% Question 14
\question
Figure 3 shows a vessel in the form of a frustum which is made by cutting off the lower part of an inverted right circular cone of base radius $72\text{ cm}$ and height $96\text{ cm}$. The height of the vessel is $60\text{ cm}$. The vessel is placed on a horizontal table. Some water is now poured into the vessel. John finds that the depth of water in the vessel is $28\text{ cm}$.

\begin{center}
\begin{tikzpicture}[scale=0.85]
% Top ellipse (z = 96, r = 3.6)
\draw[thick] (0, 3.6) ellipse (3.6cm and 0.65cm);
% Water surface ellipse (z = 64, r = 2.4)
\draw[thick] (-2.4, 2.4) arc (180:360:2.4cm and 0.43cm);
\draw[thick, densely dashed] (2.4, 2.4) arc (0:180:2.4cm and 0.43cm);
% Vessel base ellipse (z = 36, r = 1.35)
\draw[thick] (-1.35, 1.35) arc (180:360:1.35cm and 0.24cm);
\draw[thick, densely dashed] (1.35, 1.35) arc (0:180:1.35cm and 0.24cm);

% Slanted sides
\draw[thick] (-3.6, 3.6) -- (-1.35, 1.35);
\draw[thick] (3.6, 3.6) -- (1.35, 1.35);
\draw[thick, densely dotted] (-1.35, 1.35) -- (0, 0) -- (1.35, 1.35);

% Center vertical line and radii
\draw[thick, densely dotted] (0, 0) -- (0, 3.6);
\draw[thick, densely dotted] (0, 3.6) -- (3.6, 3.6) node[midway, above] {\small $72\text{ cm}$};
\draw[thick, densely dotted] (0, 2.4) -- (2.4, 2.4);
\draw[thick, densely dotted] (0, 1.35) -- (1.35, 1.35);

% Right angle symbols
\draw (0, 3.35) -- (0.25, 3.35) -- (0.25, 3.6);
\draw (0, 2.18) -- (0.22, 2.18) -- (0.22, 2.4);
\draw (0, 1.18) -- (0.17, 1.18) -- (0.17, 1.35);

% Dimension lines
\draw[thick, dashed] (2.4, 2.4) -- (3.1, 2.4);
\draw[thick, dashed] (1.35, 1.35) -- (4.3, 1.35);
\draw[thick, dashed] (3.6, 3.6) -- (5.2, 3.6);
\draw[thick, dashed] (0, 0) -- (5.2, 0);

\draw[{Stealth}-{Stealth}, thick] (2.8, 1.35) -- (2.8, 2.4) node[midway, right, rotate=90, yshift=-7pt, xshift=-8pt] {\small $28\text{ cm}$};
\draw[{Stealth}-{Stealth}, thick] (3.9, 1.35) -- (3.9, 3.6) node[midway, right, rotate=90, yshift=-7pt, xshift=-8pt] {\small $60\text{ cm}$};
\draw[{Stealth}-{Stealth}, thick] (4.8, 0) -- (4.8, 3.6) node[midway, right, rotate=90, yshift=-7pt, xshift=-8pt] {\small $96\text{ cm}$};

\node[below=8pt] at (0, 0) {Figure 3};
\end{tikzpicture}
\end{center}

\begin{parts}
\part Find the area of the wet curved surface of the vessel in terms of $\pi$. \rmarks{4}
\part John claims that the volume of water in the vessel is greater than $0.1\text{ m}^3$. Do you agree? Explain your answer. \rmarks{4}
\end{parts}

\begin{solution}
\begin{parts}
\part The slant height of the circular cone is
\[
\sqrt{72^2 + 96^2} = 120\text{ cm} \quad \text{(1M)}
\]
The area of the wet curved surface of the vessel is
\begin{align*}
&= \pi(72)(120)\frac{(96 - 60 + 28)^2 - (96 - 60)^2}{96^2} \quad &\text{(1M + 1M)} \\
&= \pi(72)(120)\frac{64^2 - 36^2}{96^2} \\
&= 2625\pi\text{ cm}^2 \quad &\text{(1A)}
\end{align*}
\part The volume of the circular cone is
\[
\frac{1}{3}\pi(72)^2(96) = 165888\pi\text{ cm}^3 \quad \text{(1M)}
\]
The volume of water in the vessel is
\begin{align*}
&= 165888\pi\left(\frac{64^3 - 36^3}{96^3}\right) \quad &\text{(1M + 1A)} \\
&= 40404\pi\text{ cm}^3 \\
&\approx 0.126932909\text{ m}^3 > 0.1\text{ m}^3
\end{align*}
Thus, the claim is agreed. \hfill (1A f.t.)
\end{parts}
\end{solution}
